\documentclass[11pt,letterpaper]{article}
\usepackage[margin=1in]{geometry}
\usepackage{setspace}
\usepackage{enumitem}
\usepackage{array}
\usepackage{booktabs}
\usepackage{tabularx}
\usepackage{graphicx}
\usepackage{titlesec}
\usepackage{fancyhdr}
\usepackage[T1]{fontenc}
\usepackage[utf8]{inputenc}
\usepackage{lmodern}

\setstretch{1.03}
\pagestyle{empty}
\setlength{\parindent}{15pt}
\setlength{\parskip}{4pt}
\setlist{nosep,leftmargin=0.24in}
\titlespacing*{\section}{0pt}{0.35em}{0.1em}
\titlespacing*{\subsection}{0pt}{0.25em}{0.08em}
\titleformat{\section}{\normalfont\bfseries\large}{}{0pt}{}
\titleformat{\subsection}{\normalfont\bfseries}{}{0pt}{}

\begin{document}
\begin{center}
{\Large\bfseries Biographical Sketch}\\[0.05in]
{\normalsize Mostafa Rezaali}\\
{\normalsize NSF AGS-PRF Application}
\end{center}
\vspace{-0.08in}
\section*{Biographical Sketch Common Form}
\textbf{Name:} Mostafa Rezaali\\
\textbf{Position Title:} Graduate Research Assistant and Ph.D. Candidate\\
\textbf{Organization:} University of Florida\\
\textbf{Source basis:} Prepared from SciENcv record 2685636 visible content and local CV publication details. The SciENcv page lists the University of Florida Ph.D. preparation entry, University of Florida appointments, and an empty related-products table; products below are therefore populated from the CV for completeness.

\section*{A. Professional Preparation}
\begin{tabularx}{\textwidth}{p{0.30\textwidth}p{0.21\textwidth}p{0.23\textwidth}X}
\toprule
\textbf{Organization} & \textbf{Location} & \textbf{Degree / Training} & \textbf{Field / Date}\\
\midrule
University of Florida & Gainesville, Florida & Doctor of Philosophy & Department of Geography, Aug 2026\\
University of Florida & Gainesville, Florida & Graduate Certificate & Atmospheric Sciences\\
Qom University of Technology & Qom, Iran & M.Sc. & Civil and Environmental Engineering, 2018\\
IAUKHSH & Isfahan, Iran & B.Sc. & Civil and Environmental Engineering, 2016\\
\bottomrule
\end{tabularx}

\section*{B. Appointments and Positions}
\begin{tabularx}{\textwidth}{p{0.16\textwidth}p{0.31\textwidth}p{0.34\textwidth}X}
\toprule
\textbf{Dates} & \textbf{Title} & \textbf{Organization / Department} & \textbf{Location}\\
\midrule
2022--2026 & GRA & University of Florida & Gainesville, Florida\\
2022--Present & Graduate Research Assistant (PhD Candidate) & University of Florida, Department of Geography & Gainesville, Florida\\
2025 & NSF LEAP REU Mentor & University of Florida & Gainesville, Florida\\
2025 & Invited Lecturer, Deep Learning Applications in Climate Science & University of Florida & Gainesville, Florida\\
2025 & Invited Lecturer, Climate Change Impacts on Future Plant Distributions & University of Florida & Gainesville, Florida\\
2021--Present & M.Sc. Student Advisor & Alborz University of Medical Sciences & Iran\\
\bottomrule
\end{tabularx}

\section*{C. Products}
\textbf{Lead-authored products related to the proposed project, sorted by date.}
\begin{enumerate}
\item Rezaali, M., Fouladi-Fard, R., O'Shaughnessy, P., Naddafi, K., and Karimi, A. (2025). Assessment of AERMOD and ADMS for NOx dispersion modeling with a combination of line and point sources. \textit{Stochastic Environmental Research and Risk Assessment}.
\item Rezaali, M., Jahangir, M. S., Fouladi-Fard, R., and Keellings, D. (2024). An ensemble deep learning approach to spatiotemporal tropospheric ozone forecasting: A case study of Tehran, Iran. \textit{Urban Climate}, 55, 101950.
\item Rezaali, M., Fouladi-Fard, R., and Karimi, A. (2023). Performance of TANN, NARX, and GMDHT models for urban water demand forecasting. \textit{Avicenna Journal of Environmental Health Engineering}, 10(2), 85--97.
\item Rezaali, M., Quilty, J., and Karimi, A. (2021). Probabilistic urban water demand forecasting using wavelet-based machine learning models. \textit{Journal of Hydrology}, 600, 126358.
\item Rezaali, M., Fouladi-Fard, R., Mojarad, H., Sorooshian, A., Mahdinia, M., et al. (2021). A wavelet-based random forest approach for indoor BTEX spatiotemporal modeling and health risk assessment. \textit{Environmental Science and Pollution Research}, 28, 22522--22535.
\end{enumerate}

\textbf{Other significant products.}
\begin{enumerate}
\item Narayanan, A., Rezaali, M., Bunting, E. L., and Keellings, D. (2025). It's getting hot in here: Spatial impact of humidity on heat wave severity in the U.S. \textit{Science of The Total Environment}, 963, 178397.
\item Farajollahi, M., Fahiminia, M., Fouladi-Fard, R., Rezaali, M., and Sorooshian, A. (2024). Human and ecological health-risk research related to environmental exposure.
\item Rezaali, M., and Fouladi-Fard, R. (2021). Aerosolized SARS-CoV-2 exposure assessment: dispersion modeling with AERMOD. \textit{Journal of Environmental Health Science and Engineering}, 19, 285--293.
\item Rezaali, M., Fouladi-Fard, R., and Karimi, A. (2020). Identification of temporal and spatial patterns of river water quality parameters using NLPCA and multivariate statistical techniques. \textit{International Journal of Environmental Science and Technology}, 17, 2977--2994.
\item Rezaali, M., and Karimi, A. (2019). Decentralized wastewater treatment plants site selection of Qom Province using environmental and spatial criteria.
\end{enumerate}

\section*{D. Synergistic Activities}
\begin{enumerate}
\item Mentored emerging climate-data scientists through NSF LEAP REU activities and student research advising.
\item Delivered invited instruction on deep learning applications in climate science and climate-change impacts.
\item Developed interdisciplinary AI workflows for climate, hydrology, air quality, and environmental health.
\item Provided peer-review service for journals including \textit{Scientific Reports}, \textit{Journal of Hydrology}, and \textit{Springer Nature Applied Sciences}.
\item Built reproducible Python, MATLAB, R, and geospatial workflows for large climate and environmental datasets.
\end{enumerate}

\section*{Eligibility and Certification Context}
Mostafa Rezaali is a Ph.D. candidate applying as an early-career researcher transitioning from doctoral training to postdoctoral research and does not hold a tenure-track or equivalent permanent research position. The Ph.D. preparation entry in SciENcv lists University of Florida, Department of Geography, with expected receipt in August 2026. The final submission should confirm the exact degree timeline and use the certified SciENcv export if Research.gov requires SciENcv certification.
\end{document}
